% 哥白尼原则（综述）
% license CCBYSA3
% type Wiki

本文根据 CC-BY-SA 协议转载翻译自维基百科\href{https://en.wikipedia.org/wiki/Copernican_principle}{相关文章}。

在物理宇宙学中哥白尼原理（Copernican principle）指出：人类并非宇宙中的特权观察者，\(^\text{[1]}\)来自地球的观测结果可视为来自宇宙中平均位置的代表性观测。该原理以哥白尼的日心说（heliocentrism）命名，是一种工作假设，其思想源自对哥白尼“地球运动”论证的宇宙学延伸。\(^\text{[2]}\)
\subsection{起源与意义}
赫尔曼·邦迪（Hermann Bondi）在二十世纪中期将该原理命名为“哥白尼原理”，但该思想本身可追溯至十六至十七世纪期间从托勒密体系（Ptolemaic system）向新范式的转变——后者不再将地球置于宇宙中心。哥白尼提出，行星的运动可通过假设太阳处于静止的中心位置来解释，从而与地心说形成对比。他认为，行星的视逆行运动是一种由地球绕太阳运动所造成的幻象，而在哥白尼模型中，太阳被置于宇宙的中心。值得注意的是，哥白尼的主要动机是出于对旧体系的技术性不满，而非出于支持任何“平庸原理”（principle of mediocrity）的哲学立场。\(^\text{[3]}\)

尽管哥白尼的日心模型常被描述为“将地球降级”——使其失去在托勒密地心模型中的中心地位，但真正采纳这一新宇宙观的，是哥白尼的继承者，尤其是十六世纪的焦尔达诺·布鲁诺（Giordano Bruno）。在传统观念中，地球的中心位置被认为是“最低且最污秽之处”；而伽利略则指出，地球实际上参与了“群星的舞蹈”，而非停留在“宇宙污秽与浮渣的积聚之所”。\(^\text{[4][5]}\)

至二十世纪后期，卡尔·萨根（Carl Sagan）以诗意的语言进一步强化了哥白尼原理的精神，他问道：“我们是谁？我们发现自己居住在一颗微不足道的行星上，这颗行星环绕着一颗平凡的恒星，处于银河系的某个被遗忘的角落，而宇宙中星系的数量远远超过人类的数量。”\(^\text{[6]}\)

尽管哥白尼原理（Copernican principle）起源于对过去假设的否定——如地心说（geocentrism）、日心说（heliocentrism）或银河中心说（galactocentrism）等，这些假设都认为人类处于宇宙的中心——但哥白尼原理的内涵比非中心论（acentricism）更为深刻。非中心论仅仅主张人类不位于宇宙中心，而哥白尼原理在承认非中心论的前提下，更进一步断言：人类观察者，或来自地球的观测结果，是宇宙中“平均位置”观测的代表。迈克尔·罗文-罗宾逊（Michael Rowan-Robinson）强调，哥白尼原理是现代思想的门槛性检验标准。他指出：“显而易见，在人类历史的后哥白尼时代，没有一个见多识广且理性的人会认为地球在宇宙中占据着独一无二的位置。”\(^\text{[7]}\)

现代宇宙学的大部分理论建立在一个前提之上：宇宙学原理（cosmological principle）在最大尺度上几乎成立，但并非绝对成立。哥白尼原理代表了支持这一前提所需的、不可约的哲学假设——当其与观测结果结合时，构成了宇宙学的逻辑基础。如果假定哥白尼原理成立，并且从地球的视角观测到宇宙在各个方向上都是各向同性的（isotropic，即在各个方向上看起来相同），则可以推断宇宙在总体上是均匀的（homogeneous，即在任意给定时刻处处相似），并且围绕任意一点都是各向同性的。这两种条件共同构成了宇宙学原理。\(^\text{[7]}\)

实际上，天文学家观测到的宇宙在星系超星系团（galactic superclusters）、纤维状结构（filaments）以及巨大空洞（great voids）等尺度上表现出明显的不均匀性。在当代主流宇宙学模型——即Lambda-CDM 模型（$\Lambda$CDM model）中，宇宙被预测为在更大的观测尺度上将逐渐趋于均匀和各向同性，当观测尺度超过约 $260$ 百万秒差距（$\mathrm{Mpc}$）时，可探测到的结构将极少。\(^\text{[8]}\)然而，来自星系团、\cite{ref9,ref10} 类星体（quasars）、\(^\text{[11]}\)以及 Ia 型超新星（type Ia supernovae）\(^\text{[12]}\)的最新证据表明，宇宙在大尺度上可能违反了各向同性。此外，科学家还发现了多种大尺度结构，如 Clowes–Campusano LQG、斯隆长城（Sloan Great Wall）、\emph{U1.11}、巨大类星体群（Huge-LQG）、武仙–北冕座大墙（Hercules–Corona Borealis Great Wall）、巨弧（Giant Arc）以及\emph{本地空洞}（Local Hole），\cite{ref13,ref14,ref15,ref16} 这些结构都暗示着宇宙的均匀性可能被破坏。
